(1) 在 \(r>0\) 无源区域，径向方程为 \(\phi''+2\phi'/r-\kappa^2\phi=0\)，解是 \((Ae^{-\kappa r}+Be^{\kappa r})/r\)。无穷远衰减要求 \(B=0\)，原点小球高斯通量要求 \(4\pi\varepsilon A=Q\)，故 \(\phi(r)=Qe^{-\kappa r}/(4\pi\varepsilon r)\)。
(2) 电场径向分量 \(E_r=-\phi'=Qe^{-\kappa r}(1+\kappa r)/(4\pi\varepsilon r^2)\)。高斯定律给包围总电荷 \(Q_{\rm enc}(r)=4\pi\varepsilon r^2 E_r=Qe^{-\kappa r}(1+\kappa r)\)，其中包括原点自由电荷和内部离子响应，不能简单说始终等于裸点电荷。
(3) 屏蔽电荷密度为 \(\rho_{\rm ion}=-\varepsilon\kappa^2\phi=-Q\kappa^2e^{-\kappa r}/(4\pi r)\)。半径 \(r\) 内离子电荷 \(Q_{\rm ion}(r)=Q_{\rm enc}(r)-Q=Q[e^{-\kappa r}(1+\kappa r)-1]\)，在 \(r\to\infty\) 时趋 \(-Q\)，故点电荷与屏蔽云总电量为零。
(4) 固定 \(r\) 取 \(\kappa\to0\)，指数因子趋一，恢复 \(\phi=Q/(4\pi\varepsilon r)\)、\(E_r=Q/(4\pi\varepsilon r^2)\)。若 \(r\gg\kappa^{-1}\)，势按 \(e^{-\kappa r}/r\) 快速衰减；线性化要求局部电势能相对热能足够小（典型条件 \(|q\phi|\ll k_BT\)），靠近强点电荷时可能失效。
